\documentclass[preprint,12pt]{elsarticle}
% Shared set-up of the journal versions (Elsevier elsarticle class). Loaded after \documentclass{elsarticle}.
\usepackage{amsmath,amssymb,bm,mathtools}
\usepackage{fontspec}
\setmainfont{TeX Gyre Termes}\setsansfont{TeX Gyre Heros}\setmonofont{Latin Modern Mono}[Scale=0.9]
\usepackage{graphicx,float,caption,subcaption,booktabs,longtable,array,multirow,tabularx,rotating}
\usepackage{xcolor,setspace,lineno,microtype,xurl,siunitx,enumitem,pdflscape}
\definecolor{lk}{RGB}{0,70,140}
\usepackage[colorlinks=true,linkcolor=lk,citecolor=lk,urlcolor=lk,breaklinks=true,pdfborder={0 0 0}]{hyperref}
\usepackage[nameinlink,capitalise]{cleveref}
\crefname{figure}{Fig.}{Figs.}\Crefname{figure}{Fig.}{Figs.}
\crefname{equation}{Eq.}{Eqs.}\Crefname{equation}{Eq.}{Eqs.}
\sisetup{detect-all,separate-uncertainty=false}
\captionsetup{font=small,labelfont=bf}
\setlength{\emergencystretch}{3em}
\newcommand{\figdir}{../figures}

\newcommand{\doihl}[1]{doi: \href{https://doi.org/#1}{\nolinkurl{#1}}}
\newcommand{\dd}{\mathrm{d}}
\newcommand{\ET}{\mathit{ET}}
\newcommand{\Ent}{\mathit{ET}_{\mathrm{nt}}}
\newcommand{\Ks}{K_{\mathrm{s}}}
\newcommand{\Se}{S_{\mathrm{e}}}
\newcommand{\TAW}{\mathit{TAW}}
\newcommand{\repourl}{https://github.com/halilunaziru73-creator/Coupled-Water-Fluxes-Pore-Scale-Dynamics-Precision-Irrigation-Mulching-Feedback-Loops}
\newcommand{\repo}{\href{\repourl}{repository}}
\newcommand{\dataurl}{\repourl/tree/main/Data}
\newcommand{\codeurl}{\repourl/tree/main/Manuscripts}
\newcommand{\softurl}{\repourl/tree/main/Software}
\newcommand{\datalink}{\href{\dataurl}{Data}}
\newcommand{\codelink}{\href{\codeurl}{Code}}
\newcommand{\softlink}{\href{\softurl}{Software}}
\newcommand{\manuallink}{\href{\manualurl}{Manual}}
\newcommand{\downloadlink}{\href{\exeurl}{SWAFL.exe}}
\newcommand{\manualurl}{https://github.com/halilunaziru73-creator/Coupled-Water-Fluxes-Pore-Scale-Dynamics-Precision-Irrigation-Mulching-Feedback-Loops/raw/main/Software/v1.1.0/SWAFL_Desktop/docs/manual_v1.1.0/SWAFL_Manual_v1.1.0.pdf}
\newcommand{\exeurl}{https://github.com/halilunaziru73-creator/Coupled-Water-Fluxes-Pore-Scale-Dynamics-Precision-Irrigation-Mulching-Feedback-Loops/releases/latest/download/SWAFL.exe}
\newcommand{\csurl}{https://github.com/halilunaziru73-creator/Coupled-Water-Fluxes-Pore-Scale-Dynamics-Precision-Irrigation-Mulching-Feedback-Loops/tree/main/Data/Comprehensive_Summary}
\newcommand{\cslink}{\href{\csurl}{\nolinkurl{Data/Comprehensive_Summary}}}
\newcommand{\csavail}{\par\smallskip\textbf{Full simulation outputs.} Every value of every log of every run (the annual water and nitrogen balances, the harvests and the crop state, year by year), one compressed table per experiment (\texttt{E00\_Baselines} to \texttt{E16\_Edge\_Cases}), the one-row-per-run summary of all 10\,190 runs (\texttt{00\_summary\_all\_simulations}), the per-experiment summary (\texttt{00\_SUMMARY\_ALL\_EXPERIMENTS}) and the script that made them (\texttt{summarize\_results3.py}) are in the folder \cslink{} of the repository (\url{\csurl}); their 20-year means reproduce the run table within 2\,\%.}

% the page number at the top right of every page, the first page too
\makeatletter
\def\ps@plain{\def\@oddhead{\hfil{\small\thepage}}\let\@evenhead\@oddhead\let\@oddfoot\@empty\let\@evenfoot\@empty}
\g@addto@macro\ps@pprintTitle{\def\@oddhead{\hfil{\small\thepage}}\let\@evenhead\@oddhead}
\pagestyle{plain}
\makeatother

\journal{SoftwareX}
\hypersetup{pdftitle={SWAFL Daisy Workbench: a graphical user interface that plans, runs and explains Daisy soil-plant-atmosphere simulations},pdfauthor={Naziru Halilu}}
\begin{document}

\begin{frontmatter}
\title{SWAFL Daisy Workbench: a graphical user interface that plans, runs and explains Daisy soil--plant--atmosphere simulations}

\author[aua,upna]{Naziru Halilu\corref{cor1}}
\ead{halilunaziru73@gmail.com}
\cortext[cor1]{Corresponding author. ORCID: 0009-0006-2324-9599.}
\affiliation[aua]{organization={Agricultural University of Athens},city={Athens},country={Greece}}
\affiliation[upna]{organization={Public University of Navarre},city={Pamplona},country={Spain}}

\begin{abstract}
Process models such as Daisy are specified in text files, run from a command line and read from log files, which keeps them out of reach of most people with a question about a field. The SWAFL Daisy Workbench is an open, offline graphical interface that runs the Daisy model itself. A question is turned into a design that is checked against a catalogue of what Daisy can vary and kept within a run budget set by the user; Daisy is run and every file it writes is kept; and the outputs are read back into a summary of the effect of every input on every result, journal-style figures of every logged quantity, and the root cause of every failed run. A bridge hands the runs to decision layers under one rule: a Daisy result is shown only from a Daisy run.
\end{abstract}

\begin{keyword}
Daisy model \sep graphical user interface \sep simulation experiments \sep model coupling \sep decision support \sep software verification
\end{keyword}
\end{frontmatter}

\section*{Code metadata}
\begin{table}[H]\centering\small
\begin{tabularx}{\textwidth}{lp{5.2cm}X}
\toprule
Nr. & Code metadata description & Metadata\\\midrule
C1 & Current code version & v1.1.0\\
C2 & Permanent link to code/repository used for this code version & \url{\repourl}\\
C3 & Permanent link to Reproducible Capsule & Not available; checks in \texttt{Manuscripts/reproducibility/p456}\\
C4 & Legal Code License & MIT (workbench); Daisy is distributed under its own licence\\
C5 & Code versioning system used & git\\
C6 & Software code languages, tools, and services used & JavaScript (Electron, Node.js), HTML/CSS; Python for the checks; Daisy 7.1.14\\
C7 & Compilation requirements, operating environments \& dependencies & Node.js and Electron (build); Daisy 7.1.14 for Windows is bundled; runs in a web browser for everything except running Daisy\\
C8 & If available, link to developer documentation/manual & \url{\manualurl}\\
C9 & Support email for questions & halilunaziru73@gmail.com\\\bottomrule
\end{tabularx}
\end{table}

\section*{Software metadata}
\begin{table}[H]\centering\small
\begin{tabularx}{\textwidth}{lp{5.2cm}X}
\toprule
Nr. & Software metadata description & Metadata\\\midrule
S1 & Current software version & 1.1.0\\
S2 & Permanent link to executables of this version & \url{\exeurl}\\
S3 & Permanent link to Reproducible Capsule & Not available\\
S4 & Legal Software License & MIT (source); end-user licence for the compiled application\\
S5 & Computing platforms/Operating Systems & Microsoft Windows (portable executable with Daisy); web browser (planning, reading and plotting of runs)\\
S6 & Installation requirements \& dependencies & None for the Windows executable (no installation, no administrator rights, no network)\\
S7 & If available, link to user manual & \url{\manualurl}\\
S8 & Support email for questions & halilunaziru73@gmail.com\\\bottomrule
\end{tabularx}
\end{table}

\linenumbers
\section{Motivation and significance}
Process-based models of the soil--plant--atmosphere system are the main tools for asking what would happen to the water, the crop and the nitrogen of a field if management or climate changed. Daisy couples water and heat flow in the soil, soil organic matter and nitrogen, the crop and the surface, and has been described, calibrated and tested in a long series of studies \cite{Hansen1990,Abrahamsen2000,Hansen2012,Holbak2026}; comparable models serve crops and irrigation \cite{Jones2003,Keating2003,Steduto2009} and water and solute flow \cite{Simunek2008,Simunek2016}. In all of them a simulation is specified in a text file, run by a program and read from files. That is natural for a model developer and a barrier for a person with a question about a field: how much of the water a mulch saves is lost to drainage, which trigger suits a sandy soil, what a smaller nitrogen dose does to nitrate leaching.

The barrier has a second part. A set of simulations that answers a question (a factorial design over trigger, sensor depth and rain, say) is a piece of work that is rarely recorded with the answer. A published ensemble of 10\,190 Daisy runs in 17 experiments \cite{HaliluData} makes such designs explicit, but a reader who wants to repeat one experiment with another soil or weather has no easy way to do so. And when a model result reaches a decision tool by hand-made steps, its origin is easily lost, or a number from another source is shown under the name of the model.

The SWAFL Daisy Workbench addresses these problems. It is part of the open SWAFL Desktop application and is designed around four principles: questions before syntax (a user states a question and the set-up file is the result, not the starting point); everything that runs is shown first (a design lists its runs, settings and estimated time before anything runs); one source of data (site weather and soil are entered once); and no hidden steps between a result and a decision.

\section{Software description}
\subsection{Software architecture}
The workbench is a window of SWAFL Desktop (\cref{fig:system}) organised in six groups of tabs (\cref{fig:wb}). A guided path of nine steps runs from the weather and soil of the site, through the library of soils, crops and fertilisers, scenarios, experiments, the process map, log plots, mass balance and run comparison, to the hand-over to the decision layers. A second row holds the tools of model development: an editor for Daisy set-up (\texttt{.dai}) files, a one-dimensional Richards soil column, response curves, calibration, and documentation of 107 reference topics. The Windows program is a single portable executable that contains the Electron runtime, the application, Daisy 7.1.14 with its library (about 47\,MB, run from its own folder) and the manual, and needs no installation, administrator rights or network. Seven weather records from five countries are bundled, and users can add their own weather and soil.

\begin{figure}[!htbp]\centering
\includegraphics[width=\textwidth]{\figdir/p6/fig01_system.pdf}
\caption{Architecture of SWAFL Desktop 1.1.0. The field tools (green) call the engines and write to the project; the project (blue bar) carries the state of the Properties panel and the weather to the five computing parts (brown); these feed the layers (blue); the layers and the workbench shell talk through one bridge; the Daisy workbench (dashed frame) hands its runs to the Daisy layers.}\label{fig:system}
\end{figure}
\begin{figure}[!htbp]\centering
\includegraphics[width=\textwidth]{\figdir/p4/fig01_workbench.pdf}
\caption{Modules of the Daisy workbench and the hand-over. The chain runs from the planner and the decision base through the experiments, Daisy and the parsers to the hand-over; the modules below the dotted line are used before Daisy is started.}\label{fig:wb}
\end{figure}

The \emph{bridge} (\cref{fig:bridge}) has five parts: a parser of Daisy logs (tab-separated \texttt{.dlf} files); a seasonal summariser that reduces the logs of a run to one record per season (harvest, water balance terms, nitrogen and chemical budgets); a soil map that links the soil of a Daisy run to the soils of the decision layers; a question map that lists, for five kinds of question, the layers that answer them; and a live object holding the summarised runs, which is the only source from which the Daisy layers read. This enforces the rule of origin: a layer shows a Daisy result only from a Daisy run.

\begin{figure}[!htbp]\centering
\includegraphics[width=\textwidth]{\figdir/p5/fig01_bridge.pdf}
\caption{The bridge from a Daisy run to a decision. Heavy arrows are the path of a run, dashed arrows are outputs and the dashed arrow at the left is the return path of the weather of the project to Daisy.}\label{fig:bridge}
\end{figure}

\subsection{Software functionalities}
\begin{itemize}[nosep]
\item \textbf{Planning}: scenarios from the Daisy library (manures, pesticides, PFAS, crops, soils), 14 prepared experiments of the published ensemble as compact designs of 6 to 72 runs, and designs of the user's own; every factor and level is checked against a catalogue of 54 factors of what Daisy can vary, and a plain-words planner proposes the factors, station and soil of a question.
\item \textbf{Run budget}: the number of runs set by the user is a ceiling that includes the rainfed references; when the combinations are more, a seeded balanced (Latin-hypercube), random, spread or first-combinations choice is taken, so a plan always gives the same runs.
\item \textbf{Soil-hydraulic sensitivity}: five factors scale the van~Genuchten--Mualem parameters of every horizon ($K_\mathrm{sat}$, $\alpha$, $n-1$, $\theta_s$ and a shift of $\ell$) by editing only the hydraulic blocks of the soil file.
\item \textbf{Runs kept}: every run is a folder of its own with its set-up, the place where every value was written and the files Daisy wrote, next to the Daisy that ran it.
\item \textbf{Outputs}: a summary of the range of every result and the effect of every input on every result, a ranking of combinations by a weighted score, journal-style figures of every quantity of every log (by factor level, through the year, monthly, yearly, cumulative, depth--time maps), and mass-balance checks.
\item \textbf{Diagnosis}: a run without results is diagnosed from the Daisy message; the offending line of the set-up is located and shown, the cause is classed (unknown name, solver failure, value out of range, missing file, weather, duplicate, unreadable line, silent stop), and runs failing for the same cause are grouped with the change that removes it.
\item \textbf{Model-development tools}: the set-up editor, the Richards soil column, the Daisy denitrification response curves, and calibration by sequential search, CMA-ES \cite{HansenOstermeier2001} or Bayesian optimisation \cite{Mockus1978}.
\end{itemize}

\section{Illustrative examples}
\subsection{Running the published experiments on a laptop}
The study behind the ensemble took 17.5 days of Daisy time (median 145\,s per run). The 14 compact designs of the workbench (\cref{tab:designs}) contain 368 runs that stand for 7\,990 runs of the study, about 15 hours instead of 361 hours. Rebuilding the marginal response curves of five outputs from the compact levels recovered them with a median error of 13.2\,\% of their range over 51 combinations of experiment, factor and response, smallest where the design has enough levels (sensor depth, 4.9\,\%) and largest where three levels must describe a bending response (dose, 30\,\%): the compact designs are for orientation, and a curve that matters is re-run with more levels where it bends, which is one setting of the workbench. \Cref{fig:wbsteps} shows the guided path on a real case.

\begin{table}[!htbp]\centering\small
\caption{The 17 experiments of the study and their prepared designs in the workbench. Hours are the sum of the Daisy wall-clock run times of the study (\texttt{run\_seconds}) and the estimate for the compact design (runs $\times$ the median run time of the experiment). The experiments marked $-$ are not offered as designs; their published runs are bundled.}\label{tab:designs}
\begin{tabular}{llrrrr}\toprule
Experiment & Question & Study runs & Study (h) & Compact runs & Compact (h)\\\midrule
E0 & Baselines & 10 & 0.6 & 9 & 0.6\\
E1 & Dose and duration & 200 & 11.0 & 17 & 0.9\\
E2 & Trigger pressure & 200 & 7.7 & 21 & 0.8\\
E3 & Sensor depth & 200 & 8.1 & 21 & 0.8\\
E4 & Methods & 100 & 4.1 & 6 & 0.2\\
E5 & Soil structure dynamics (SHP) & 300 & 15.8 & 72 & 3.5\\
E6 & Mediterranean pollution & 210 & 6.9 & 30 & 0.8\\
E7 & Mulching & 200 & 8.9 & $-$ & $-$\\
E8 & Rain intensity and moisture & 1000 & 1.7 & $-$ & $-$\\
E9 & Soil moisture and rain & 1000 & 2.0 & 40 & 0.1\\
E10 & Amount, timing and soil & 1000 & 70.4 & 20 & 0.7\\
E11 & Temperature, evaporation, rain & 1000 & 42.0 & 16 & 0.7\\
E12 & Sensor, trigger, rain & 1000 & 42.9 & 30 & 1.2\\
E13 & N, water, soil & 1000 & 57.4 & 32 & 1.9\\
E14 & Mulch and rain & 1000 & 48.0 & $-$ & $-$\\
E15 & Method, depth, amount & 1000 & 43.6 & 30 & 1.3\\
E16 & Edge cases & 770 & 48.2 & 24 & 1.2\\
\midrule Total & & 10190 & 419 & 368 & 14.7\\\bottomrule\end{tabular}\end{table}
\begin{figure}[!htbp]\centering
\begin{subfigure}{0.49\textwidth}\centering\includegraphics[width=\textwidth]{\figdir/shots_jpg/wb_step1_input.jpg}\caption{Step 1: my data and weather (manual p.~760).}\end{subfigure}\hfill
\begin{subfigure}{0.49\textwidth}\centering\includegraphics[width=\textwidth]{\figdir/shots_jpg/wb_step2_library.jpg}\caption{Step 2: library of soils, crops and fertilisers (p.~583).}\end{subfigure}\\[3pt]
\begin{subfigure}{0.49\textwidth}\centering\includegraphics[width=\textwidth]{\figdir/shots_jpg/wb_step4_scenarios.jpg}\caption{Step 4: scenarios, a question and its runs (p.~762).}\end{subfigure}\hfill
\begin{subfigure}{0.49\textwidth}\centering\includegraphics[width=\textwidth]{\figdir/shots_jpg/wb_step6_logplots.jpg}\caption{Step 6: log plots of the runs (p.~763).}\end{subfigure}\\[3pt]
\begin{subfigure}{0.49\textwidth}\centering\includegraphics[width=\textwidth]{\figdir/shots_jpg/wb_step9_layers.jpg}\caption{Step 9: to the SWAFL layers, which layer answers which question (p.~764).}\end{subfigure}\hfill
\begin{subfigure}{0.49\textwidth}\centering\includegraphics[width=\textwidth]{\figdir/shots_jpg/wb_tool_setup_editor.jpg}\caption{Tool: set-up editor with bracket check (p.~783).}\end{subfigure}\\[3pt]
\begin{subfigure}{0.49\textwidth}\centering\includegraphics[width=\textwidth]{\figdir/shots_jpg/wb_tool_calibrate.jpg}\caption{Tool: calibration (daisypy-optim) (p.~779).}\end{subfigure}\hfill
\begin{subfigure}{0.49\textwidth}\centering\includegraphics[width=\textwidth]{\figdir/shots_jpg/wb_tool_learn_files.jpg}\caption{Tool: learn the Daisy files, the grammar of a .dai file (p.~792).}\end{subfigure}
\caption{The guided path of the Daisy workbench and three of its tools (real screenshots of release 1.1.0). Each step carries its step number, a one-line ``what to do here'', the manual page and Back and Next buttons.}\label{fig:wbsteps}
\end{figure}

Daisy 7.1.14 ran the set-ups written by the workbench for the library scenarios, deficit irrigation and the soil-hydraulic factors, and each finished and wrote its logs. Scaling the hydraulic parameters of a loam towards a faster-draining soil raised matrix percolation of the first season from 177 to 222\,mm, removed runoff and lowered actual evapotranspiration, the expected direction. A defect that stopped every run on one horizon of the Daisy samples was found and removed.

\subsection{Reading and explaining the runs}
For a designed experiment the workbench produces the summary of effects, the figures of every log and the diagnosis of failed runs (\cref{fig:engine}). The bridge reads the logs exactly: six bundled logs were read by the bridge and independently by \texttt{pandas}, with relative differences of column sums of at most $3\times10^{-15}$; seasonal nitrogen summaries matched the log columns to their rounding, and the chemical budget closed within 0.02\,\% of the applied amount. Two bundled nitrogen runs with the same fertiliser and almost the same crop uptake differed by a factor of 3.9 in leaching (\cref{tab:nrun}), which a comparison of uptake alone would miss.

\begin{figure}[!htbp]\centering
\begin{subfigure}{\textwidth}\centering\includegraphics[width=0.92\textwidth]{\figdir/shots_jpg/wb_summary_effects.jpg}\caption{The summary of an experiment: the effect of every input on every result (blue rises, red falls; in brackets the share of the mean) and the summary in words.}\end{subfigure}

\vspace{4pt}
\begin{subfigure}{\textwidth}\centering\includegraphics[width=0.92\textwidth]{\figdir/shots_jpg/wb_every_log_lai.jpg}\caption{One of the figures of every log: leaf area index through the year, one panel per water limit, rainfed and irrigated, the mean of the seasons with 10 to 90\,\% as the band.}\end{subfigure}
\caption{Outputs of the experiment engine drawn from Daisy runs (screen captures of the workbench).}\label{fig:engine}
\end{figure}
\begin{table}[!htbp]\centering\small
\caption{Mean seasonal nitrogen budget of the two bundled annual nitrogen logs as reported by the bridge (kg\,N\,ha$^{-1}$, ten seasons).}\label{tab:nrun}
\begin{tabular}{lrrrrrr}\toprule
Run & Fertiliser & Mineralisation & Deposition & Crop uptake & Denitrification & Leaching\\\midrule
3b & 229.0 & 136.0 & 13.2 & 341.5 & 15.8 & 4.4\\
4b & 229.0 & 134.9 & 13.2 & 338.5 & 14.5 & 17.0\\\bottomrule
\end{tabular}
\end{table}

\subsection{Checks against independent references}
The soil column conserved mass (daily error below 0.001\,mm) and differed from an independent method-of-lines reference by 0.3\,\% in drainage for sand and 0.9\,\% for loam after six days (\cref{fig:richards}); for a ponded silty clay loam it does not converge with node spacing and shows its mass-balance error so that this is visible. The denitrification response curve was compared with 792 inputs and outputs stored by the Daisy project: the first comparison showed a median error of 40\,\% caused by rescaling inputs that Daisy had already rescaled; after correction all 792 rows matched to $3\times10^{-15}$. Calibration recovered a well-constrained parameter to 0.0002 to 0.03 decades and showed, through its search history, that two poorly constrained parameters were not identifiable from the stored outputs. The planner returned the expected factors, station and soil for 8 of 8 test questions.

\begin{figure}[!htbp]\centering
\includegraphics[width=\textwidth]{\figdir/p4/fig04_richards.pdf}
\caption{Soil column. (a) Drainage below 1\,m with the node spacing; dashed lines are the reference. (b) Final storage minus the reference. (c) Infiltration of 200\,mm of rain into a ponded silty clay loam (left axis) and run time (right axis), showing the lack of convergence.}\label{fig:richards}
\end{figure}

\section{Impact}
The workbench makes Daisy usable by people who are not model developers: a question becomes a checked design, the runs are kept and explained, and a published experiment can be repeated on another soil or weather in hours on a laptop. Three routes from Daisy to a decision become available in one program (\cref{tab:routes}): the published ensemble through a meta-model in microseconds per scenario, compact experiments in hours, and the user's own runs in minutes each. This supports research questions that need designed simulation experiments (sensitivity to hydraulic parameters, deficit irrigation, nitrogen and chemical fate), teaching of soil--plant--atmosphere modelling, where a student sees the set-up, the run and the result together, and model development, because the diagnosis of failed runs and the comparison with stored outputs locate errors that are invisible on the screen, as the denitrification defect showed. The bridge design (log parser, seasonal summariser, soil map, question map and a rule of origin) is not specific to Daisy and can be reused for other process models. Limits are stated: the workbench was tested in a browser with real Daisy outputs and Daisy was run end to end for selected scenarios, but not for every factor; Daisy is bundled for Windows only; and no field measurements were used.

\begin{table}[!htbp]\centering\small
\caption{The three routes from Daisy to a decision.}\label{tab:routes}
\begin{tabularx}{\textwidth}{>{\raggedright\arraybackslash}p{3.0cm}XXX}
\toprule & Ensemble and Engine 1 & Compact experiment (workbench) & Own live run\\\midrule
Soil and weather & ten soils, one weather record, the design of the study & as chosen in the workbench (library or own soil, bundled or own weather) & own soil table and weather\\
Time & microseconds per scenario & about 0.2 to 4 hours per experiment & minutes per run\\
Needs Daisy & no & yes (Windows) & yes (Windows)\\
Answers & ranking of options, response curves, regime maps within the design & the shape of a response for the chosen conditions & the budget of one set-up, for the layers of the suite\\
Main limit & fixed design; offset of absolute levels & coarse levels (workbench) & one run, one set-up\\\bottomrule
\end{tabularx}
\end{table}

\section{Conclusions}
The SWAFL Daisy Workbench is an open, offline graphical interface that plans, runs and explains Daisy simulations. It checks every design against a catalogue of what Daisy can vary, keeps an exact run budget, keeps every run file by file, and returns a summary of effects, journal-style figures of every logged quantity and the root cause of failures. A bridge hands the runs to decision layers under the rule that a Daisy result is shown only from a Daisy run. Each computing part was verified against an independent reference or stored Daisy output, and the verification found and removed a defect in the shared denitrification function. Repeating the published experiments run by run, testing the installer on more Windows systems, and comparing runs with field measurements are the next steps.

\section*{CRediT authorship contribution statement}
\textbf{Naziru Halilu:} Conceptualization, Methodology, Software, Validation, Formal analysis, Data curation, Visualization, Writing -- original draft, Writing -- review \& editing.

\section*{Declaration of competing interest}
The author declares that he has no known competing financial interests or personal relationships that could have appeared to influence the work reported in this paper.

\section*{Declaration of generative AI and AI-assisted technologies in the writing process}
During the preparation of this work the author used Claude (Anthropic) to edit language and to format the manuscript for the journal. After using this tool, the author reviewed and edited the content as needed and takes full responsibility for the content of the publication.

\section*{Funding}
This research did not receive any specific grant from funding agencies in the public, commercial, or not-for-profit sectors.

\section*{Data availability}
The published simulation ensemble is archived on Zenodo (\doihl{10.5281/zenodo.23043626}); every log of every run is in the repository folder \cslink{}. The checks of this article are in \url{\codeurl}.

\section*{Acknowledgements}
The Daisy model is developed and maintained at the University of Copenhagen.

\bibliographystyle{elsarticle-num}
\bibliography{../tex/refs}
\end{document}
